\section{Discussion}

This section interprets the empirical findings with respect to the research
questions and the existing literature. The limitations of the study are then
discussed, followed by implications for future research and potential
application scenarios.

\subsection{Interpretation of Classification Performance and Explainability}

RQ1 examined how the investigated ML classifiers differ in performance for
NLP-based prioritization of IT support tickets. On the common data, split, and
feature basis, LightGBM achieved the highest hold-out Macro-F1 on FS4 and was
the only alternative that significantly outperformed Logistic Regression after
Holm correction. The observed difference was
$\Delta F1_{\mathrm{macro}} = 0.0197$, which supports H1. This does not
indicate a general superiority of more complex model families. Earlier ticket
studies likewise report context-dependent rankings, with different models
performing best across datasets and classical SVM approaches remaining
competitive with more complex methods
\cite{ahmed2023incident,martins2024port,bonechi2024support}.

The feature basis also affected performance. FS4 achieved the highest average
cross-validation performance across all seven classifiers, indicating that
queue information added predictive value in this dataset. The SHAP analysis
of the final model also assigned a relevant share of attribution mass to the
queue. This is consistent with work that incorporates ticket-related and
contextual information in addition to text
\cite{beckers2009multifacet,fuchs2022support}. Queue structures are tied to
organizational processes, so this predictive value may not transfer directly
to other service organizations. Sentiment features added little.
FS5-LightGBM reached a hold-out Macro-F1 of 0.5597 compared with 0.5594 for
FS4-LightGBM, while sentiment accounted for less than one percent of absolute
SHAP mass.

The final model should not be interpreted as ready for immediate production
use. Low-priority tickets were identified less reliably, with a Recall of
0.3291, and Macro-F1 does not reflect different operational consequences of
misclassification. The moderate classification performance may partly reflect the information
available for priority prediction, as \citet{tian2015priority} model
bug-report priority as a multi-factor decision involving textual, temporal,
author, related-report, severity, and product information. In contrast, the
present study was restricted to the text and structured features available in
the investigated dataset. In addition, real-world ITSM research reports
particular difficulties for semantically complex and underrepresented classes
\cite{razma2026classification}. RQ1 is therefore answered by the finding that the
classifiers differed in prioritization performance. LightGBM performed best
on the common feature basis and was the only alternative with a statistically
significant advantage over Logistic Regression, while no general superiority
of more complex model families was observed.

RQ2 examined differences in SHAP-based explanation properties. Explanation
compactness was operationalized for inferential testing. Logistic Regression
showed significantly higher Top-10 concentration than Complement Naive Bayes
and Random Forest, which supports H2 under the prespecified criterion. No such
advantage was observed over Linear SVM, kNN, XGBoost, or LightGBM.
Robustness analyses confirmed the result for Complement Naive Bayes across all
five operationalizations, whereas the comparison with Random Forest depended
on the compactness measure used.

These findings extend earlier XAI work in ticket classification. Local
explanation methods have been used to analyze relevant features and
misclassifications, while SHAP-based explanations have been integrated into
helpdesk systems \cite{zicari2022ticket,licina2026helpagent}. In contrast to studies that primarily evaluate ticket-classification performance
or use explanations within individual classification frameworks \cite{razma2026classification,zicari2022ticket,licina2026helpagent}, the present
analysis compares SHAP-based explanation compactness across the investigated
classifiers. This comparison shows that predictive performance and explanation compactness
do not follow a consistent trade-off. Explanation quality is multidimensional
\cite{nauta2023evaluation}, so faithfulness and robustness were considered in
addition to compactness. The faithfulness analysis indicated that highlighted
features were related to model behavior, while local cases showed
heterogeneous plausibility. SHAP describes feature attributions to model predictions
\cite{lundberg2017shap}, which should not be interpreted as establishing
causal relationships \cite{janzing2020causal}. Comprehensibility for
actual users was not examined systematically, so the findings do not replace
independent user or expert evaluation. RQ2 is therefore answered by
model-dependent differences in SHAP-based explanation properties without a
consistent advantage of Logistic Regression or a simple separation between
linear and more complex models.

\subsection{Limitations}

A central limitation is the synthetic data basis. Public availability supports
reproducible comparison \cite{bueck2025support}. Synthetic data may not fully
capture characteristics and distributions observed in real-world data
\cite{goyal2024synthetic}. In the present context, this concerns in particular
language, support processes, and error patterns and may therefore limit
transferability to real ITSM environments. The dataset also includes broader queues such as
Human Resources, Sales and Pre-Sales, and General Inquiry. The findings
therefore concern a heterogeneous support population rather than a narrowly
defined IT service desk.

The analysis assumed that the queue was already available when priority was
predicted. Because no process timestamps establish the order of queue and
priority assignment, this assumption could not be verified. Priority and
language labels could not be validated independently. Manual reviews assessed
the plausibility of data-quality decisions rather than reassigning priorities.
Semantic grouping also depends on the Sentence-BERT representation, the
similarity threshold of 0.95, and the complete-linkage strategy
\cite{reimers2019sentencebert}.

The group-based hold-out measures generalization within the same dataset.
Neither external nor temporal validation was conducted. This is relevant because
model performance in IT-related classification has been shown to depend on the
application context \cite{ahmed2023incident}, while temporal changes may alter
the relationship between input data and target concepts \cite{gama2014conceptdrift}.
Consequently, concept drift, alternative queue structures, and
organization-specific prioritization rules remain unobserved. The algorithm and representation screenings were exploratory and budget-limited. Hyperparameters selected on FS1 were fixed for later feature-set comparisons. This improves comparability but may prevent individual
models from reaching their optimum on later feature sets. The cross-validation
Macro-F1 of 0.5330 therefore represents model-selection performance within the
training partition, while the hold-out Macro-F1 of 0.5597 represents the
separate final evaluation.

Different model structures required different SHAP explainers, which limits
the direct comparability of the H2 results across models. Because
Shapley-based attributions can vary depending on how they are computed
\cite{sundararajan2020many}, the observed H2 differences in explanation
compactness may reflect both model differences and differences in the
attribution method. The H2 sample contained 100 semantic
groups per priority class, giving all classes equal weight rather than their
natural frequencies. Random Forest and XGBoost were explained approximately
for computational reasons. Exact and approximate XGBoost attributions showed
high agreement with $\rho = 0.9666$, but this cannot be assumed for Random
Forest. The kNN explanations were sensitive to the background sample and
evaluation budget. The qualitative plausibility assessment was conducted by the author rather
than independent ITSM experts. Since explanation quality is multidimensional
\cite{nauta2023evaluation}, the present analysis cannot establish whether the
explanations are understandable or decision-relevant for ITSM practitioners.

\subsection{Implications and Future Research}

The findings suggest that automated ticket prioritization can benefit from
structured contextual information in addition to text. More complex
classifiers did not consistently outperform the baseline, so a
supportive human-in-the-loop setting appears more plausible initially than
fully autonomous priority assignment. This approach was also adopted in prior work
on intelligent ticket management \cite{zicari2022ticket}.

Future studies should test the findings on real-world ITSM data across multiple
organizations and use temporal test sets. Cost-sensitive or ordinal learning
objectives could better reflect unequal operational consequences of errors.
The exploratory transformer screening could also be extended by evaluating
multilingual transformer classifiers such as XLM-RoBERTa (XLM-R) \cite{conneau2020xlmr}
under the full group-based modeling and hold-out design. Independent assessments
by ITSM experts would complement model-based XAI metrics by examining
comprehensibility and actual decision utility.

Overall, the study provides a controlled comparison of classifiers and their SHAP-based explanations under a common evaluation design. Whether the observed technical differences translate into operational value requires subsequent real-world process and user evaluation.